\section{TVM: la apertura de la compilación de aprendizaje automático}

\subsection{La motivación que parte de un problema}

De CXXNet a MXNet, Chen Tianqi (陈天奇) descubrió un cuello de botella que se repetía: cada generación de GPU y de hardware nuevo requería una gran cantidad de kernels escritos a mano, que debían admitir distintos layouts de datos (como NHWC frente a NCHW), distintos tamaños de entrada y distintas formas de hardware; el espacio de combinaciones era enorme. Los kernels escritos a mano «sencillamente no bastaban».

\begin{importantbox}{El problema central que TVM debía resolver}
¿Se podía usar la técnica de compilación para generar automáticamente estos kernels? ¿Se podía lograr que menos personas construyeran rápidamente, con la técnica correcta, sistemas de aprendizaje automático orientados a hardware nuevo? Este era el punto de partida de TVM------\textbf{que la complejidad de ingeniería de los sistemas de aprendizaje automático dejara de crecer linealmente con la cantidad de tipos de hardware}.
\end{importantbox}

\subsection{La diferencia entre GPU, TPU y NPU, y su convergencia}

En la entrevista, Chen Tianqi explicó la diferencia entre tres tipos de chips de cómputo:

\begin{itemize}
    \item \textbf{GPU}: proviene de la programación gráfica; su característica central es SIMT (Single Instruction Multiple Thread), que programa con una gran cantidad de hilos en paralelo. Con el desarrollo de CUDA, ya «no tiene relación» con la programación gráfica.
    \item \textbf{TPU/NPU}: cuenta con un conjunto de instrucciones más complejo, donde una sola instrucción puede completar una operación matricial masiva o una lectura/escritura de memoria a gran escala. Por lo general no tiene el concepto de hilo, porque una sola instrucción ya es, de por sí, muy «grande».
    \item \textbf{Tendencia de convergencia}: las GPU incorporaron el Tensor Core (que originalmente solo tenían las NPU), y las NPU también empezaron a incorporar el concepto de hilo para acercarse al ecosistema de CUDA. «Ahora este límite es cada vez más difuso.»
\end{itemize}

Esta tendencia de convergencia es precisamente la manifestación del valor persistente de TVM------el modelo de programación de cada generación de hardware cambia, y se necesita la técnica de compilación para tender un puente entre las diferencias.

\subsection{El proceso de desarrollo: «construir una ciudad en una isla desierta»}

TVM se inspiró en parte en Halide (un trabajo de Jonathan Ragan-Kelley y otros, del MIT, que usaba la técnica de compilación para optimizar pipelines en el campo del procesamiento de imágenes), pero la compilación de aprendizaje automático era, en ese momento, un terreno completamente vacío. Chen Tianqi describió esta experiencia como «construir una ciudad en una isla desierta»------no había ningún camino ya trazado al que recurrir.

Desde que empezó a idearlo, al volver del GTC de abril de 2016, hasta la primera versión ejecutable, en marzo de 2017, pasaron cerca de 11 meses, la mayor parte dedicados a escribir código. Recuerda que este proceso fue de «iteración constante»: «Aunque no lo entendiera, no tenía miedo; era como un becerro recién nacido que no le teme al tigre.»

Mientras hacía TVM, junto con Thierry Moreau, su colaborador del área de arquitectura de computadoras, construyó con FPGA un acelerador de NPU programable. Su idea inicial era más «wild»------meter directamente la red neuronal, mediante bake, en el cableado (wire) del FPGA (más adelante alguien hizo un trabajo relacionado con HLS). Al final decidieron hacer primero una NPU parecida a una TPU, pero con una filosofía de diseño completamente distinta de la corriente dominante en ese momento.

\begin{warningbox}{El pago de la experiencia del trabajo grande de la clase ACM}
Al recordar el desarrollo de TVM, Chen Tianqi mencionó varias veces su experiencia en la clase ACM. «En ese entonces, en la clase ACM teníamos un trabajo grande de arquitectura de computadoras, y yo me propuse el reto de ejecutar, en un CPU que yo mismo había escrito, un compilador que yo mismo había escrito. TVM se parece muchísimo a esa experiencia------también era construir desde cero un compilador for un hardware nuevo. Gracias a esa experiencia anterior, sentí que no me ponía tan nervioso.» Esto muestra que un proyecto retador que, en la etapa de licenciatura, parece una «actividad extracurricular», puede, años después, dar un pago de una forma inesperada.
\end{warningbox}

Chen Tianqi señaló en particular que, en ese entonces, él «no venía formado en arquitectura de computadoras»; pensaba por completo from first principle cómo hacerlo. El conductor comentó: «Siempre has sido así------en el concurso de bachillerato eras un becerro recién nacido que no le teme al tigre, escribiendo el compilador en la clase ACM eras un becerro recién nacido que no le teme al tigre, haciendo XGBoost eras un becerro recién nacido que no le teme al tigre, haciendo TVM sigues siendo un becerro recién nacido que no le teme al tigre. ¿No es esta tu mayor ventaja?» Chen Tianqi respondió: «Creo que elegir el problema y cómo abordarlo es muy importante. Hacer un problema mediocre y hacer un problema importante toman casi el mismo tiempo, así que hay que elegir sin falta el problema correcto.»

Al mismo tiempo, junto con su colaborador del área de arquitectura de computadoras, construyó con FPGA un acelerador de NPU programable (VTA). El planteamiento de diseño de esta NPU tenía una visión muy adelantada a su tiempo:

\begin{knowledgebox}{El diseño innovador del acelerador VTA}
La mayoría de las NPU tempranas eran completamente hardwired (no programables, con una sola función fija). Pero VTA adoptó un diseño \textbf{programable a nivel de instrucción}, que permitía implementar múltiples operadores con un mismo conjunto de microinstrucciones. Para resolver el problema del traslape (overlap) entre la lectura/escritura de memoria y el cómputo, se inspiraron en el concepto de hyper-threading: dividieron lógicamente la NPU en dos hilos, uno dedicado a las operaciones de memoria y otro al cómputo, sincronizados mediante un mecanismo de señales de hardware. Este diseño «sigue siendo bastante interesante incluso hoy», y en la programación asíncrona más reciente de NVIDIA se pueden ver rastros similares.
\end{knowledgebox}

\subsection{El sistema de coordenadas de Jeff Dean y la elección de un tema de alto riesgo}

Durante su estancia como pasante en Google Brain, Jeff Dean llevó a Chen Tianqi a una sala de juntas y dibujó un sistema de coordenadas: un eje era el riesgo (risk), el otro el impacto (impact). Entre las distintas opciones de proyecto, Chen Tianqi eligió el tema con mayor riesgo y también mayor impacto------Net2Net. El problema de este proyecto era: dado un modelo pequeño ya entrenado, ¿cómo hacer un bootstrap rápido de un modelo grande, de modo que su capacidad pudiera aumentar?

La razón por la que Chen Tianqi eligió este tema fue que siempre tuvo una visión: la inteligencia necesita poder evolucionar, y la capacidad del modelo necesita poder crecer. Si una red neuronal empieza siendo pequeña, en cierto punto deja de ser suficiente. Desde la perspectiva del lifelong learning, ¿cómo resolver este problema? Señaló: «De hecho, ahora te das cuenta de que este problema todavía existe de la misma manera------esta es también la razón por la que ahora hay una tendencia hacia el scaling up a partir de modelos pequeños.»

Jeff Dean era, en ese entonces, su mentor en Brain, y además «tenía la mente muy abierta». Chen Tianqi fue, en ese momento, prácticamente el primer intern de Jeff Dean------en una etapa en la que «su GAN tal vez todavía no se había publicado, o apenas se acababa de publicar».

Su criterio de elección era: \textbf{«El tiempo es limitado, así que hay que hacer sin falta lo que uno quiere hacer. No tiene que ser algo con impacto, pero sí tiene que ser algo que uno quiera hacer.»}

Ante la pregunta de «si les recomendarías a los investigadores jóvenes elegir temas de alto riesgo y alto rendimiento», su respuesta fue muy pragmática: puede ser distinto para cada persona. Pero como él ya había pasado por muchos fracasos, «fracasar una vez más tampoco es tan grave». «Almost always terminamos eligiendo una dirección de riesgo relativamente alto y retorno relativamente alto.»

\subsection{El valor durante el posdoctorado y la historia de la Recomputation}

Antes de hacer TVM, Chen Tianqi había pensado en un método para «entrenar modelos más profundos con una complejidad menor que lineal» (ahorrando memoria mediante recomputation), pero Carlos le dijo que este trabajo, publicado tal cual, «sería apenas un NeurIPS poster», que no era suficientemente profundo, y que si se iba a hacer, había que llevarlo al nivel de best paper------había que diseñar un método más automatizado para elegir la estrategia de recomputation. Como su energía era limitada, no continuó por esa línea, y más adelante hubo un follow-up work que hizo un trabajo relacionado.

La exigencia de Carlos no era negar el valor de esa línea, sino decir: \textbf{«Tienes que elegir lo que haces, y una vez que decides hacerlo, tienes que hacerlo a fondo.»}

Li Mu (李沐), en un comentario en Zhihu, reveló que, en el tercer o cuarto año de su doctorado, Chen Tianqi le dijo que iba a hacer TVM, y Li Mu le advirtió que el riesgo era demasiado alto------abrir un pozo nuevo y grande cuando el doctorado estaba por terminar. La respuesta de Chen Tianqi fue: «Sin importar si en el futuro voy a trabajar, a emprender o a buscar una posición académica, lo más importante sigue siendo hacer, en ese momento, lo que uno mismo considera más importante.» Li Mu comentó después: «Visto ahora, esta es la única forma correcta de actuar.»

\begin{importantbox}{La lógica de fondo para tomar decisiones}
Al tomar decisiones, Chen Tianqi se hace una pregunta a sí mismo: \textbf{«Si elijo este camino y fracaso, ¿me voy a arrepentir?»} Si la respuesta es que no se arrepentiría, lo elige. Esto no es un análisis racional de costo-beneficio, sino una manera de asegurarse de seguir lo que siente en el corazón. «Muchas veces de lo que uno se arrepiente no es de si logró el éxito, sino de si tomó o no el camino que idealmente quería tomar.»
\end{importantbox}

\subsection{La evolución de TVM y su nicho ecológico}

Desde su nacimiento hasta ahora, TVM ha pasado por múltiples rondas de reestructuración, desde la compilación de grafos inicial hasta Relax (un compilador de grafos rediseñado para los LLM). Chen Tianqi enfatizó que la filosofía de diseño de TVM es \textbf{«evolucionable»}------reinventarse continuamente para adaptarse a las nuevas necesidades. Señaló: «Si comparas el TVM actual con su forma en el momento en que empezó, ya hemos pasado por varias generaciones de iteración.»

Frente a competidores como Torch Compile, XLA y Triton, considera que la necesidad de la técnica de compilación no solo no ha disminuido, sino que va en aumento:

\begin{itemize}
    \item \textbf{Especialización (specialized) del hardware}: la forma en que avanza cada generación de chips es distinta, y el modelo de programación cambia incluso entre generaciones sucesivas de un mismo vendor
    \item \textbf{El «entrelazamiento» entre el modelo y el hardware}: para alcanzar la eficiencia óptima, el diseño del modelo se acopla cada vez más con las características del hardware------DeepSeek, al diseñar en función de un entorno de hardware específico, es un ejemplo típico de esto
    \item \textbf{Hardware Lottery}: el hardware limita las posibilidades de desarrollo del modelo, y se necesita la técnica de compilación para romper esta restricción
    \item \textbf{El auge de los Agentes de IA}: la aparición de los sistemas de agentes y de los DSL trae nuevos escenarios de aplicación para la técnica de compilación
\end{itemize}

TVM tiene actualmente dos direcciones de desarrollo: una es integrarse mejor con ecosistemas como PyTorch (abstrayendo por separado una especificación general de llamadas a funciones, para facilitar las llamadas entre distintos frameworks); la otra es lograr avances en dominios verticales (como compilar y desplegar LLM en páginas web y en dispositivos vehiculares). Su meta última es: \textbf{«hacer el desarrollo de ingeniería de aprendizaje automático tan simple como escribir una red neuronal»}------algo que, por ahora, todavía está lejos de lograrse.

\begin{warningbox}{Reinventarse a uno mismo frente a hacer patch a un sistema existente}
Dentro de OctoML, Chen Tianqi también enfrentó la disyuntiva entre reinvent y patch. Para dar soporte a los LLM, echaron abajo la parte de compilación de grafos original de TVM y rehicieron Relax. Esta decisión «no la iba a aprobar todo el mundo», porque la técnica ya existente podía resolver el problema del momento a corto plazo. Pero él insiste en que: «If you don't reinvent yourself, there will be a time where you bite the technical debt.» A largo plazo, la acumulación de deuda técnica es más temible que el costo a corto plazo de una reestructuración.
\end{warningbox}

\subsection{Resumen de la sección}

TVM es la obra que reúne, de la manera más completa, el estilo de investigación de Chen Tianqi de «partir del problema». Surgió de un cuello de botella de ingeniería descubierto durante el desarrollo de MXNet, se inspiró en el campo de los compiladores, y en su diseño fusionó la idea del codiseño de hardware y software. Aunque enfrenta una competencia intensa, el lugar central de la técnica de compilación en los sistemas de aprendizaje automático ya es un consenso.
